% SPDX-License-Identifier: Apache-2.0
\section{Products at Their Operands' Widths}
\label{sec:x-smul}

A product of extended operands, \code{a.sext::<64>().mul::<64>(b.sext::<64>())}
for two 32-bit values, is a 32 by 32 signed product, and the netlist
says so: \code{\$signed(a) * \$signed(b)} in Verilog and
\code{signed(a) * signed(b)} in VHDL, resized to the product's width.
Written as the product of the extended words, a 64 by 64 product,
it hid the operands' widths from Vivado, which built it of four DSP
blocks in series (issue 1649). A zero-extended operand beside a signed
one is signed with a zero above it, so it stays positive. A product
narrower than its two narrow values together is written as before,
since it is not their whole product.

The unit registers three products: two signed words, a signed word and
an unsigned half, and a signed byte and an unsigned byte into sixteen
bits. The run asserts each against the
same arithmetic on Rust's integers, for operands chosen at the ends of
their ranges, and the netlist is checked against the run under both
simulators.

\begin{figure}[htbp]
\centering
\input{smul_timing.tex}
\caption{Two signed words, an unsigned half, a signed and an unsigned
byte, and the products registered from them.}
\label{fig:x-smul}
\end{figure}

\lstinputlisting[caption={\code{ex\_smul.rs}. Run with \code{bazel run //lib/examples:ex\_smul}.},label={lst:x-smul}]{lib/examples/ex_smul.rs}

\lstinputlisting[style=ascii,caption={What \code{ex\_smul} printed when the build ran it: each case and its three products, then the lowering.},label={lst:x-smul-out}]{out_smul.txt}
